%------------------------------------------------------------------------------%
\begin{question}
  Dado o vetor $\vec{u} = \overrightarrow{PQ}$ com
  \(
    P = ( 2, -1 )
  \)
  q
  \(
    Q = ( 5, 3 )
  \)
  determine

  1) o vetor $\vec{u}$

  2) o ponto na extremidade do vetor $\vec{u}$ partindo da origem

  3) o ponto na extremidade do vetor $\vec{u}$ partindo do ponto $( -1, -2)$
\end{question}

%------------------------------------------------------------------------------%
\begin{resolution}{Solução}
  \[
    \vec{u}
    \pause
    \;=\; \overrightarrow{AB}
    \pause
    \;=\; \vetor{5 - 2}{3 - (-1)}
    \pause
    \;=\; \vetor{3}{4}
  \]

  \pause
  O vetor $\vec{u}$ partindo da origem termina no ponto $(3, 4)$

  \pause
  O vetor $\vec{u}$ partindo do ponto $(-1, -2)$ termina no ponto com coordenadas
  \[
    \pause
    x
    \pause
    \;=\; -1 + u_x
    \pause
    \;=\; -1 + 3
    \pause
    \;=\; 2
  \]
  \[
    \pause
    y
    \pause
    \;=\; -2 + u_y
    \pause
    \;=\; -2 + 4
    \pause
    \;=\; 2
  \]
\end{resolution}

%------------------------------------------------------------------------------%
\endinput
